\documentclass[12pt]{article}
\usepackage[utf8]{inputenc}
\usepackage[T1]{fontenc}
\usepackage{hyperref}
\title{Carolina Algonquian Language Resources for K-12 Curriculum Development}
\author{Compiled by Claude Sonnet 5.5 (claude-sonnet-5-5) for Andr\'as Kornai}
\date{}
\begin{document}
\maketitle

\section{Introduction}
Carolina Algonquian (ISO 639-3 crr; Glottocode caro1243) is an extinct Eastern Algonquian language formerly spoken in coastal North Carolina (Roanoke, Croatoan, Secotan and related groups). It is attested only in short word lists, chiefly those of Thomas Harriot from 1585--86. The search results revealed no living speakers and no revitalization programme.

\section{Primary Educational Resources}
No primers, textbooks or school programmes were found. Any heritage curriculum would have to rest on the small attested vocabulary.

\section{Linguistic and Reference Materials}
\begin{itemize}
\item Encyclopedia overview of the language: \url{https://en.wikipedia.org/wiki/Carolina_Algonquian_language}
\item Article on Thomas Harriot and the lost North Carolina Algonquian language: \url{https://brewminate.com/thomas-harriot-and-the-lost-north-carolina-algonquian-language/}
\item Article ``Thomas Harriot (1560--1621) and the English origins of Algonkian linguistics'' (John Benjamins): \url{https://www.benjamins.com/online/hl/articles/hl.19.1.03sal}
\end{itemize}
The search summary also mentions Scott Dawson's 2014 paper ``The Vocabulary of Croatoan Algonquian'' (120 words and phrases) but no URL for it was seen.

\section{Religious and Community Resources}
\begin{itemize}
\item David S. Phelps, ``The Carolina Algonkians'' (historical essay): \url{https://homepages.rootsweb.com/~jmack/algonqin/phelps1.htm}
\end{itemize}
No community language organisations were found.

\section{Digital and Audio Resources}
No audio, software or corpus resources were found.

\section{Curriculum Development Context}
With no speakers and only fragmentary records, a K-12 programme would be a heritage and history unit with a small word list rather than a full language curriculum. Reconstruction from related Eastern Algonquian languages would require specialist work not covered by these results.

\end{document}